\documentclass[10pt,journal,compsoc]{IEEEtran}

\usepackage[utf8]{inputenc}
\usepackage[spanish,es-tabla]{babel}
\usepackage{amsmath,amsfonts,amssymb}
\usepackage{graphicx}
\usepackage{booktabs}
\usepackage{tabularx}
\usepackage{multirow}
\usepackage{xcolor}
\usepackage{caption}
\usepackage{subcaption}
\usepackage{microtype}
\usepackage[colorlinks=true,linkcolor=blue,citecolor=blue,urlcolor=blue]{hyperref}

\title{Compendio de Evidencia Experimental, Tablas y Figuras Científicas: Ateneo+}
\author{Jorge Doicela, et al.\\\textit{Ateneo+: Simulador Clínico Multimodal con RAG Híbrido, Motor KST/BKT y Analítica del Aprendizaje}}

\begin{document}

\maketitle

\begin{abstract}
Este documento compila la totalidad de las tablas empíricas y figuras de alta resolución (300 DPI) generadas por el pipeline automatizado de MLOps y analítica de Ateneo+. Los resultados sustentan la evaluación de rendimiento del recuperador denso fine-tuned (\texttt{ateneo-bge-m3-ecuador-v2}) sobre las 42 Guías de Práctica Clínica (GPC) oficiales del Ministerio de Salud Pública (MSP) del Ecuador indexadas en 7,052 fragmentos normativos indivisibles, la auditoría de fidelidad normativa (\textit{Faithfulness Score}), la ganancia de aprendizaje psicométrica (\textit{Hake Learning Gain}) y la modelación probabilística mediante \textit{Knowledge Space Theory} (KST) y \textit{Bayesian Knowledge Tracing} (BKT).
\end{abstract}

\section{Fine-Tuning Supervisado y Evaluación del Recuperador Denso (BGE-M3 Ecuador v2)}

El modelo recuperador \texttt{ateneo-bge-m3-ecuador-v2} fue optimizado en una GPU NVIDIA A100 (80GB VRAM) mediante la función de pérdida \textit{Multiple Negatives Ranking Loss} (MNRL, $\tau=0.02$) a partir de un banco canónico de 1,710 pares/tripletas clínicas supervisadas con minería de \textit{hard negatives} intra e inter-especialidad, alcanzando una concordancia inter-anotador de $\kappa_w = 0.9509$. La Tabla~\ref{tab:pre_post_retrieval} documenta la comparación empírica antes y después del fine-tuning en el conjunto de prueba ciego independiente.

% TABLA I (PRE vs POST FINE-TUNING)
\begin{table}[htbp]
\centering
\caption{Comparación empírica Pre vs. Post Fine-Tuning en Ateneo+ BGE-M3 Ecuador v2 (MNRL $\tau=0.02$).}
\label{tab:pre_post_retrieval}
\begin{tabular}{lcccc}
\toprule
\textbf{Métrica de Retrieval} & \textbf{Baseline (Pre)} & \textbf{Ateneo+ (Post)} & \textbf{Delta ($\Delta$)} & \textbf{Ganancia (\%)} \\
\midrule
    Accuracy@1 (Hit@1) & 0.0522 & 0.2450 & +0.1928 & +369.23\% \\
    Accuracy@5 (Hit@5) & 0.1606 & 0.4739 & +0.3133 & +195.00\% \\
    MRR@1 & 0.0522 & 0.2450 & +0.1928 & +369.23\% \\
    MRR@5 & 0.0889 & 0.3328 & +0.2439 & +274.40\% \\
    NDCG@5 & 0.1066 & 0.3682 & +0.2616 & +245.40\% \\
    Precision@5 & 0.0321 & 0.0948 & +0.0627 & +195.00\% \\
    Recall@5 & 0.1606 & 0.4739 & +0.3133 & +195.00\% \\
\bottomrule
\end{tabular}
\vspace{1mm}
\begin{minipage}{\linewidth}
\footnotesize
\textit{Nota:} Evaluación cuantitativa estandarizada sobre el pool de validación/prueba v2 (570 tripletas). El fine-tuning contrastivo multiplica por 4.7x la capacidad de recuperación en Top-1 y triplica el Mean Reciprocal Rank (MRR@5).
\end{minipage}
\end{table}

\section{Benchmark Ciego Out-of-Distribution (283 Consultas Clínicas)}

Para evaluar la capacidad de generalización en casos no vistos durante el entrenamiento, la Tabla~\ref{tab:blind_retrieval_benchmark} reporta el rendimiento en el conjunto de prueba ciego independiente ($N=283$ consultas clínicas). La prueba de rangos con signo de Wilcoxon confirma una superioridad estadísticamente significativa de Ateneo-BGE-M3 frente al modelo base ($W=365.0, p=0.0001018$).

% TABLA II (BENCHMARK CIEGO OOD)
\begin{table}[htbp]
\centering
\caption{Evaluación Comparativa de Recuperación Normativa en el Test Set Ciego OOD (283 consultas clínicas)}
\label{tab:blind_retrieval_benchmark}
\begin{tabular}{lccccc}
\toprule
\textbf{Arquitectura} & \textbf{Hit@1} & \textbf{Hit@3} & \textbf{Hit@5} & \textbf{MRR@5} & \textbf{NDCG@5} \\
\midrule
BM25 Puro (Léxico) & 2.47\% & 4.95\% & 5.30\% & 3.55\% & 3.99\% \\
BAAI/bge-m3 (Zero-Shot) & 0.71\% & 3.18\% & 3.53\% & 1.90\% & 2.31\% \\
\textbf{Ateneo-BGE-M3 (FT)} & \textbf{4.24\%} & \textbf{9.19\%} & \textbf{11.66\%} & \textbf{7.12\%} & \textbf{8.26\%} \\
Ensamble RRF ($k=60$) & 3.53\% & 7.77\% & 9.89\% & 5.87\% & 6.88\% \\
\bottomrule
\end{tabular}
\vspace{1mm}
\begin{minipage}{\linewidth}
\footnotesize
\textit{Nota:} Ateneo-BGE-M3 supera en 6x en Hit@1 y 3.75x en MRR@5 al modelo base (Wilcoxon $p < 0.001$).
\end{minipage}
\end{table}

\section{Auditoría de Grounding Normativo (Faithfulness Score)}

La Tabla~\ref{tab:faithfulness_results} contrasta la tasa de anclaje normativo y mitigación de alucinaciones diagnósticas frente al baseline de GPT-4o Zero-Shot y un RAG genérico no especializado.

% TABLA II (FAITHFULNESS)
\begin{table}[htbp]
\centering
\caption{Auditoría de Grounding Normativo y Fidelidad de Retroalimentación RAG frente a las GPCs del MSP}
\label{tab:faithfulness_results}
\begin{tabular}{lcccc}
\toprule
\textbf{Arquitectura Evaluativa} & \textbf{Fidelidad} & \textbf{Afirmaciones} & \textbf{Alucinación} & \textbf{Seguridad} \\
\midrule
Baseline Zero-Shot (GPT-4o) & 54.2\% & 26 / 48 & 45.8\% & Riesgo Moderado \\
RAG Genérico (Embeddings Base) & 82.5\% & 38 / 46 & 17.5\% & Grounding Parcial \\
\textbf{Ateneo+ RAG Híbrido FT} & \textbf{100.0\%} & \textbf{36 / 36} & \textbf{$< 5.0\%$} & \textbf{Alto Grounding} \\
\bottomrule
\end{tabular}
\vspace{1mm}
\begin{minipage}{\linewidth}
\footnotesize
\textit{Nota:} El \textit{Faithfulness Score} evalúa la proporción de aciertos y omisiones con correlación textual directa en los fragmentos del MSP, mitigando errores diagnósticos en el entrenamiento formativo.
\end{minipage}
\end{table}

\section{Estudio Piloto de Ganancia de Aprendizaje (Learning Gain)}

La Tabla~\ref{tab:learning_gain_pilot} y la Figura~\ref{fig:learning_gain} documentan el impacto formativo medido en la cohorte piloto mediante el estadístico normalizado de Hake ($g$).

% TABLA III (PILOT STUDY)
\begin{table}[htbp]
\centering
\caption{Evaluación Cuantitativa de Ganancia de Razonamiento Clínico (Pre-Test vs. Post-Test tras Intervención con Ateneo+)}
\label{tab:learning_gain_pilot}
\begin{tabular}{lcccc}
\toprule
\textbf{Métrica Psicométrica} & \textbf{Pre-Test} & \textbf{Post-Test} & \textbf{Delta ($\Delta$)} & \textbf{Significancia ($p$)} \\
\midrule
Puntaje Global (0--10) & 4.87 $\pm$ 0.54 & \textbf{8.64 $\pm$ 0.40} & +3.77 & $p < 0.0001$ \\
Ganancia de Hake ($g$) & --- & \textbf{0.7400 $\pm$ 0.0542} & --- & \textbf{Ganancia Alta} \\
Estadístico $t$ Pareado ($df=24$) & --- & $t = 105.266$ & --- & $p < 10^{-10}$ \\
Muestra de Internos ($N$) & 25 internos & 25 internos & --- & Medicina \\
\bottomrule
\end{tabular}
\vspace{1mm}
\begin{minipage}{\linewidth}
\footnotesize
\textit{Nota:} La ganancia de Hake clasifica en \textbf{Ganancia Alta} ($g \ge 0.70$), respaldando la hipótesis de que la retroalimentación socrática y la correlación multimodal incrementan el juicio clínico estructurado.
\end{minipage}
\end{table}

\begin{figure}[htbp]
\centering
\includegraphics[width=\linewidth]{figura_learning_gain.png}
\caption{Distribución de ganancia de aprendizaje normalizada de Hake ($g$) pre-test vs. post-test tras la intervención clínica con Ateneo+.}
\label{fig:learning_gain}
\end{figure}

\section{Analítica de Aprendizaje e Índice de Brecha Formativa (IBF)}

La Figura~\ref{fig:ibf_cohorte} ilustra la evolución longitudinal del Índice de Brecha Formativa (IBF) por eje clínico en la cohorte analizada.

\begin{figure}[htbp]
\centering
\includegraphics[width=\linewidth]{figura_ibf_cohorte.png}
\caption{Evolución longitudinal del Índice de Brecha Formativa (IBF) en los 4 ejes clínicos normativos: Diagnóstico, Tratamiento, Prevención y Seguimiento.}
\label{fig:ibf_cohorte}
\end{figure}

\section{Motor de Currículo Adaptativo KST \& BKT}

La Tabla~\ref{tab:bkt_competencias} y la Figura~\ref{fig:kst_trajectory} presentan la evolución probabilística de dominio clínico comparando una ruta fija tradicional contra la ruta adaptativa en la Zona de Desarrollo Próximo (ZDP).

% TABLA IV (KST & BKT)
\begin{table}[htbp]
\centering
\caption{Probabilidad de Dominio BKT por Competencia Clínica — Ruta Fija vs. Ruta KST Adaptativa}
\label{tab:bkt_competencias}
\begin{tabular}{lcccccr}
\toprule
\textbf{Competencia} & \textbf{$L_0$} & \textbf{$P$(Fija)} & \textbf{Nivel} & \textbf{$P$(KST)} & \textbf{Nivel} & \textbf{$\Delta$KST} \\
\midrule
Semiología & 0.40 & 0.978 & \textbf{Dominado} & \textbf{0.847} & \textbf{Dominado} & -0.131 \\
Dx Diferencial & 0.30 & 0.987 & \textbf{Dominado} & \textbf{0.990} & \textbf{Dominado} & \textbf{+0.003} \\
Exámenes & 0.25 & 0.725 & En Progreso & \textbf{0.961} & \textbf{Dominado} & \textbf{+0.236} \\
Correlación & 0.15 & 0.273 & Sin Iniciar & \textbf{0.990} & \textbf{Dominado} & \textbf{+0.717} \\
Dx Final & 0.30 & 0.962 & \textbf{Dominado} & \textbf{0.990} & \textbf{Dominado} & \textbf{+0.028} \\
Tratamiento & 0.20 & 0.644 & En Progreso & \textbf{0.990} & \textbf{Dominado} & \textbf{+0.346} \\
Seguimiento & 0.25 & 0.678 & En Progreso & \textbf{0.990} & \textbf{Dominado} & \textbf{+0.312} \\
\bottomrule
\end{tabular}
\vspace{1mm}
\begin{minipage}{\linewidth}
\footnotesize
\textit{Nota:} $L_0$ = probabilidad a priori de dominio. Umbral de dominio: $P \ge 0.75$. Simulación de 10 sesiones consecutivas en ZDP (Vygotsky, 1978).
\end{minipage}
\end{table}

\begin{figure}[htbp]
\centering
\includegraphics[width=\linewidth]{figura_kst_trajectory.png}
\caption{Trayectorias de aprendizaje simuladas bajo Knowledge Space Theory (KST) y Bayesian Knowledge Tracing (BKT) demostrando la convergencia proactiva en ZDP.}
\label{fig:kst_trajectory}
\end{figure}

\end{document}
